\section{Problem Statement and the Safety Ladder}
\label{sec:method}

\subsection{Task}

DeFi workflow synthesis maps a natural-language intent $p$ to a workflow
specification $W = (N, E, C)$: a set $N$ of typed nodes drawn from a fixed
executable DeFi node vocabulary (e.g.\ \texttt{tokenSelector},
\texttt{oneInchQuote}, \texttt{priceImpact\-Calculator},
\texttt{oneInchSwap}, \texttt{limitOrder}, \texttt{fusionPlus},
\texttt{transactionMonitor}), a set $E \subseteq N \times N$ of type-level
data-flow edges, and an execution configuration $C$ mapping keys such as
\texttt{fromToken}, \texttt{amount}, \texttt{slippage}, and
\texttt{targetPrice} to concrete values.

Each benchmark item pairs a prompt with an expert-authored gold annotation
$(N^{\star}, E^{\star}, C^{\star}, S^{\star})$, where $S^{\star}$ is a set
of \emph{safety predicates} the workflow must satisfy to be safely
executable: a slippage bound, a price-impact gate, transaction monitoring,
a distinct bridge destination. The annotation also marks underspecified
prompts, for which the correct response is a clarification request rather
than a workflow.

\subsection{The safety ladder}
\label{sec:ladder}

We score every generated workflow on four strictly nested levels. Each
level subsumes the ones below it, so a system's position on the ladder
identifies which capability it lacks.

\paragraph{Level 1: graph-valid.} The workflow recalls every required node
type and required type-level edge, and the generator did not error. This
level is deliberately lenient: extra nodes are ignored unless they break a
required edge, matching the weak notion of ``valid'' that graph-only
evaluations reward.

\paragraph{Level 2: executable.} The workflow is graph-valid \emph{and}
every required configuration key in $C^{\star}$ holds a concrete,
non-placeholder value. A node that exists only in template mode, with no
trade amount or target price, is not executable.

\paragraph{Level 3: statically safe.} The workflow is executable
\emph{and} every required safety predicate in $S^{\star}$ is declared and
parameterized. Predicates are derived uniformly from the normalized
workflow for all systems, so no system benefits from self-reported safety.

\paragraph{Level 4: execution-safe.} Declaring a predicate does not prove
it gates execution, so the top level is measured by running the workflow.
We deploy a constant-product AMM (Uniswap-V2-style $x \cdot y = k$ with a
$0.3\%$ fee) and ERC20 tokens on an in-process EVM, seed deterministic
reserves, and submit each generated swap as a real transaction. The
harness observes the mined receipt, gas, or revert and labels the outcome
(\texttt{executed\_safe}, \texttt{reverted\_slippage},
\texttt{aborted\_price\_impact}, \texttt{unsafe\_executed},
\texttt{pending\_not\_filled}, or \texttt{not\_executable}). Execution
semantics are real (\texttt{transferFrom} calls, \texttt{require}
reverts, gas accounting); only the reserves are synthetic, a choice we
validate against mainnet in Section~\ref{sec:fidelity}.

A separate \emph{clarification axis} scores underspecified prompts, where
emitting a confidently wrong workflow is the failure mode.

\paragraph{Slippage is not price impact.} The distinction that drives our
headline result is between two protections that generators routinely
conflate. A \emph{slippage bound} sets a minimum output computed from the
already impact-adjusted quote; it protects against price movement between
quote and execution. It does nothing to stop a large order from executing
at severe \emph{self-inflicted} price impact. Only a separate
\emph{price-impact gate} does. A workflow that mines a high-impact trade
with no such gate is labeled \texttt{unsafe\_executed}.
